\documentclass[10pt,a4paper]{article}
\usepackage{amsmath,amssymb,amsthm,graphicx}

\usepackage{amsmath,amssymb,amsfonts,amsthm,mathtools,stmaryrd,xspace,enumitem}
\newenvironment{authorclaim}[1]{\par\medskip\noindent\textbf{#1.}\ \ignorespaces}{\par\medskip}
\newcommand{\oneset}[1]{\left\{\mathinner{#1}\right\}}
\newcommand{\ccinterval}[2]{[{#1}..{#2}]}
\newcommand{\abs}[1]{\left|\mathinner{#1}\right|}
\newcommand{\N}{\ensuremath{\mathbb N}}
\newcommand{\eps}{\varepsilon}
\newcommand{\Oh}{\mathcal O}
\newcommand{\Nc}[1]{\ensuremath{\mathsf{NC}^{#1}}\xspace}
\newcommand{\Ac}[1]{\ensuremath{\mathsf{AC}^{#1}}\xspace}
\newcommand{\DLINTIME}{\ensuremath{\mathsf{DLINTIME}}\xspace}
\newcommand{\WP}{\textsc{WP}\xspace}
\newcommand{\DTIME}{\ensuremath{\mathsf{DTIME}}\xspace}
\newcommand{\nand}{\ensuremath{\mathsf{nand}}}
\newbox\pairbox
\def\pair<#1>{{\mathsurround=0pt\setbox\pairbox\hbox{$\left\langle#1\right\rangle$}\left\langle\kern-0.35\ht\pairbox\copy\pairbox\kern-0.35\ht\pairbox\right\rangle}}
\AtBeginDocument{\author{\parbox{.95\textwidth}{\centering Laurent Bartholdi, Michael Figelius\\Markus Lohrey, Armin Wei\ss}}}
\allowdisplaybreaks
\newtheorem{definition}{Definition}

\newcommand{\Z}{\ensuremath{\mathbb Z}}
\newcommand{\supp}{\operatorname{supp}}
\newcommand{\Sym}[1]{\mathrm{Sym}({#1})}
\newcommand{\cointerval}[2]{[{#1}..{#2}-1]}
\newcommand{\LEAF}{\mathsf{LEAF}}
\newcommand{\bLEAF}{\mathsf{bLEAF}}
\newcommand{\leaf}{\mathrm{leaf}}
\newcommand{\val}{\mathrm{val}}
\newcommand{\CVP}{\textsc{CVP}\xspace}
\renewcommand{\L}{\ensuremath{\mathsf{LOGSPACE}}\xspace}
\newcommand{\NP}{\ensuremath{\mathsf{NP}}\xspace}
\newcommand{\Ptime}{\mathsf P}
\title{Wreath-product circuit value classification by universal leaf languages and the center quotient}
\author{Laurent Bartholdi, Michael Figelius, Markus Lohrey, Armin Weiß}
\date{doi:10.4230/LIPIcs.CCC.2020.29; CC BY 3.0 Unported}
\begin{document}
\maketitle
\noindent\textit{Editorial scope.} Independent formal Theorem6 universal leaf-language characterization, full AppendixC support chain grouped once. Does not claim full PSPACE corollary or uniform SENS padding proof.\par\medskip
\noindent\textit{Allowed standard prerequisites (editorial).}\begin{itemize}
\item Straight-line programs, leaf languages, standard super-decreasing subsetsum SLP encoding. Full local conventions and original Lemmas28–33 retained. Standard cited results[29,35,38] are imported with exact stated conditions; new circuit-to-subsetsum and all wreath/center witness constructions fully retained. Full51 original references.
\end{itemize}\medskip
\section*{Wreath-product circuit-value classification (Theorem 6)}
Let $G$ be a f.g.\ non-trivial group with center $Z=Z(G)$.
\begin{itemize}
\item $\mathrm{CVP}(G\wr\mathbb Z)$ belongs to $\forall\mathrm{LEAF}(\mathrm{WP}(G))$.
\item $\mathrm{CVP}(G\wr\mathbb Z)$ is hard for the class $\forall\mathrm{LEAF}(\mathrm{WP}(G/Z))$.
\end{itemize}
In particular, if $Z=1$, then $\mathrm{CVP}(G\wr\mathbb Z)$ is complete for $\forall\mathrm{LEAF}(\mathrm{WP}(G))$.

\section*{Wreath-product conventions (Section 2)} \label{sec-wreath}

A fundamental group construction that we shall use is the \emph{wreath product}: 
given groups $G$ and $H$ acting on the right on sets $X$ and $Y$ respectively, their
\emph{wreath product} $G\wr H$ is a group acting on $X\times Y$. 
We start with the restricted direct product $G^{(Y)}$ (the base group) of all mappings $f : Y \to G$ having finite support
$\supp(f) = \{ y \mid f(y) \neq 1 \}$ with the operation of pointwise multiplication. 
The group $H$ has a natural left action on
$G^{(Y)}$: for $f \in G^{(Y)}$ and $h \in H$, we define ${}^h\!f \in G^{(Y)}$ by $({}^h\!f)(y) = f(y^h)$. 
The corresponding semidirect product $G^{(Y)} \rtimes H$ is the \emph{wreath product} $G \wr H$.  In other words:
\begin{itemize}
\item
Elements of $G \wr H$ are pairs $(f,h) \in G^{(Y)} \times H$ and we simply write $fh$ for this pair.
\item
The multiplication in $G \wr H$ is defined as follows:
Let $f_1h_1, f_2h_2 \in G \wr H$. Then
$f_1h_1f_2h_2 = f_1 {}^{h_1}\!\!f_2 h_1h_2$, where
the product $f_1 {}^{h_1}\!\!f_2\colon y \mapsto f_1(y) f_2(y^{h_1})$ is the pointwise product.
\end{itemize}
The wreath product $G \wr H$ acts on $X \times Y$ by $(x,y)^{fh} = (x^{f(y)},y^h)$.
The wreath product defined above is also called the {\em (restricted) permutational wreath product}. 
There is also the variant where $G=X$ and $H=Y$ and both groups act on themselves by right-multiplication, which is called
the {\em (restricted) regular wreath product} (or {\em standard wreath product}). A subtle point is that the permutational wreath product is an associative operation whereas
the regular wreath product is in general not. The term ``restricted'' refers to the fact that the base group is $G^{(Y)}$, i.e., only finitely supported mappings are taken into account.
If $G^{(Y)}$ is replaced by $G^Y$ (i.e., the set of all mappings from $Y$ to $G$ with pointwise multiplication), then one speaks of an unrestricted wreath product.
For $Y$ finite this makes of course no difference.
We will only deal with restricted wreath products.
The action of $G$ on $X$ in the permutational wreath product is usually not important for us, but it is nice to have an associative
operation. For the right group $H$, we will only make use of the following cases: 
\begin{itemize}
\item $H = \Sym Y$ acting on $Y$,
\item $H$ a (finite or infinite) cyclic group acting on itself.
\end{itemize}
Thus, if $H$ is cyclic, the permutational wreath product and the regular wreath product (both denoted by $G \wr H$) coincide. Nevertheless, be aware that $G \wr (H \wr H) = (G \wr H) \wr H$ holds only for the permutational wreath product even if $H$ is cyclic. 
Note that if $G$ is generated by $\Sigma$ and $H$ is generated by $\Gamma$ then $G \wr H$ is 
generated by $\Sigma \cup \Gamma$.



\section*{Leaf-language conventions (Appendix A)}

In the following, we introduce basic concepts
related to leaf languages, more details can be found in  \cite{ref11,ref23,ref25,ref24,ref29}.
An NTM $M$ with input alphabet $\Gamma$ is \emph{adequate}, if 
(i) for every input $x \in \Gamma^*$, $M$ does not have an infinite 
computation on input $x$, (ii) the finite set of transition tuples
of $M$ is linearly ordered, and (iii) when terminating $M$ prints a symbol $\alpha(q)$ from a finite alphabet $\Sigma$, where $q$ is 
 the current state of $M$.
For an input $x \in \Gamma^*$, we define the computation tree
by unfolding the configuration
graph of $M$ from the initial configuration. 
By condition (i) and (ii), the computation
tree can be identified with a  
finite ordered tree $T(x) \subseteq \mathbb{N}^*$.
For $u \in T(x)$ let $q(u)$ be the $M$-state of 
the configuration that is associated with the 
tree node $u$.
Then, the leaf string $\leaf(M,x)$
is the string $\alpha(q(v_1)) \cdots \alpha(q(v_k)) \in \Sigma^+$, where 
$v_1,\ldots,v_k$ are all leaves of $T(x)$ listed in lexicographic order.

An adequate NTM $M$ is called  \emph{balanced}, if for every input $x \in \Gamma^*$,
$T(x)$ is a complete binary tree.
With a language $K \subseteq\Sigma^*$ we associate 
the language 
$$
\LEAF(M,K) = \{ x \in  \Gamma^* \mid \leaf(M,x) \in K \}.
$$
Finally, we associate two complexity classes with $K \subseteq\Sigma^*$:
\begin{eqnarray*}
\LEAF(K) &=& \{ \LEAF(M,K) \mid \text{$M$ is an adequate polynomial time
    NTM}\}\\
\bLEAF(K) &=& \{ \LEAF(M,K) \mid \text{$M$ is a balanced polynomial time NTM}\}
\end{eqnarray*}
These classes are closed under polynomial time reductions.
We clearly have $\bLEAF(K)  \subseteq \LEAF(K)$.
The following result was shown in \cite{ref29} by padding computation trees to complete binary trees.

\begin{authorclaim}{Lemma 28}
Assume that $K \subseteq \Sigma^*$ is a language such that $\Sigma$ contains a symbol $1$ with the following property:
if $uv \in K$ for $u,v \in \Sigma^*$ then $u1v \in K$. Then $\LEAF(K)  = \bLEAF(K)$. 
\end{authorclaim}
In particular, we obtain the following lemma:

\begin{authorclaim}{Lemma 29} \label{lemma-JMT}
Let $G$ be a finitely generated group and $\Sigma$ a finite standard generating set for $G$. 
Then $\LEAF(\WP(G,\Sigma))  = \bLEAF(\WP(G,\Sigma))$. 
\end{authorclaim}

Moreover, we have:
\begin{authorclaim}{Lemma 30} \label{lemma-f.g.set-not-important}
Let $G$ be finitely generated group and $\Sigma$, $\Gamma$ finite standard generating sets for $G$. 
Then $\LEAF(\WP(G,\Sigma))  = \LEAF(\WP(G,\Gamma))$. 
\end{authorclaim}

\begin{proof}
Consider a language $L \in \LEAF(\WP(G,\Sigma))$. Thus, there exists an adequate polynomial time
 NTM $M$ such that $L = \LEAF(M,\WP(G,\Sigma))$. We modify $M$ as follows: If $M$ terminates and prints
 the symbol $a \in \Sigma$, it enters a small nondeterministic subcomputation that produces the leaf string $w_a$,
 where $w_a \in \Gamma^*$ is a word that evaluates to the same group element as $a$. Let $M'$ be the resulting
 adequate polynomial time NTM.
 It follows that $\LEAF(M,\WP(G,\Sigma)) =  \LEAF(M',\WP(G,\Gamma))$.
\end{proof}
Lemma 30 allows to omit the standard generating set $\Sigma$ in the notations
$\LEAF(\WP(G,\Sigma))$ and $\bLEAF(\WP(G,\Sigma))$. We will always do that.

For a complexity class $\mathcal C$ we denote by $\forall\mathcal C$ the class of all languages $L$ such that there exists a polynomial $p(n)$ and a language $K\in\mathcal C$ with $L=\{u\mid\forall v\in\{0,1\}^{p(|u|)}:u\#v\in K\}$ (hence, for instance $\forall\mathrm P=\mathrm{coNP}$ and $\forall\mathrm{PSPACE}=\mathrm{PSPACE}$).

\section*{Straight-line programs and circuit values (Appendix B)}

Take a word $w=a_0\cdots a_{n-1}\in\Sigma^*$ over the alphabet $\Sigma$ ($n\ge0$, $a_0,\ldots,a_{n-1}\in\Sigma$). For $0\le i<n$ let $w[i]=a_i$ and for $0\le i\le j<n$ let $w[i:j]=a_ia_{i+1}\cdots a_j$. Moreover $w[:i]=w[0:i]$. Note that in the notations $w[i]$ and $w[i:j]$ we take $0$ as the first position in $w$. Finally, with $|w|_a=|\{i\mid a_i=a\}|$ we denote the number of occurrences of $a$ in $w$.
A {\em straight-line program} (SLP for short) over the alphabet $\Sigma$ is a triple $\mathcal{G} = (V,\rho,S)$, where 
$V$ is a finite set of variables such that $V \cap \Sigma = \emptyset$, $S \in V$ is the start variable, and $\rho: V \to (V \cup \Sigma)^*$ is a mapping
such that the relation $\{ (A,B) \in V \times V \colon B \text{ occurs in } \rho(A) \}$
is acyclic. For the reader familiar with context free grammars, it might be helpful to view the SLP
$\mathcal{G} = (V,\rho,S)$ as the context-free grammar $(V,\Sigma, P, S)$, where $P$ contains all productions $A \to \rho(A)$ for $A \in V$.
The definition of an SLP implies that this context-free grammar derives exactly on terminal word, which will be denoted by
$\val(\mathcal{G})$. Formally, one can extend $\rho$ to a morphism $\rho : (V \cup \Sigma)^* \to (V \cup \Sigma)^*$ by setting
$\rho(a)=a$ for all $a \in \Sigma$. The above acyclicity condition on $\rho$ implies that for $m = |V|$ we have $\rho^m(w) \in \Sigma^*$ for all
$w \in (V \cup \Sigma)^*$. We then define $\val_{\mathcal{G}}(w) = \rho^m(w)$ (the string derived from the sentential form $w$) and 
$\val(\mathcal{G}) = \val_{\mathcal{G}}(S)$.

The word $\rho(A)$ is also called the {\em right-hand side} of $A$. 
Quite often, it is convenient to assume that all right-hand sides are of the form $a \in \Sigma$ or $BC$ with $B,C \in V$.
This corresponds to the well-known Chomsky normal form for context-free grammars. There is a simple linear time 
algorithm that transforms an SLP $\mathcal{G}$ with $\val(\mathcal{G}) \neq \varepsilon$ into an SLP $\mathcal{G}'$ in Chomsky normal
form with $\val(\mathcal{G}) = \val(\mathcal{G}')$, see e.g. \cite[Proposition 3.8]{ref38}.

We define the size of the SLP $\mathcal{G} = (V,\rho,S)$ as the total length of all right-hand sides:  $|\mathcal{G}| = \sum_{A \in V} |\rho(A)|$.
SLPs offer a succinct representation of words that contain many repeated substrings. For instance, the word $(ab)^{2^n}$ can be produced
by the SLP $\mathcal{G} = (\{A_0, \ldots,A_n\},\rho,A_n)$ with $\rho(A_0) = ab$ and $\rho(A_{i+1}) = A_i A_i$ for $0 \leq i \leq n-1$.
We need the following upper bound on the length of the word $\val(\mathcal{G})$:

\begin{authorclaim}{Lemma 31 (c.f.\ [13])} \label{lemma-SLP-upper-bound}
For every SLP $\mathcal{G}$ we have $|\val(\mathcal{G})| \leq  3^{|\mathcal{G}|/3}$.
\end{authorclaim}
We also need polynomial time algorithms for a few algorithmic problems for SLPs:

\begin{authorclaim}{Lemma 32 ([38, Chapter 3])} \label{lemma-basic}
The following problems can be solved in polynomial time, where $\mathcal{G}$ is an SLP over a terminal alphabet $\Sigma$, $a \in \Sigma$, and $p, q \in \mathbb{N}$
(the latter are given in binary notation):
\begin{itemize}
\item Given $\mathcal{G}$, compute the length $|\val(\mathcal{G})|$.
\item Given $\mathcal{G}$ and $a$, compute the number $|\val(\mathcal{G})|_a$ of occurrences of $a$.
\item Given $\mathcal{G}$ and $p$, compute the symbol $\val(\mathcal{G})[p] \in \Sigma$ (in case $0 \le p < |\val(\mathcal{G})|$ does not hold, the algorithm outputs a special symbol).
\item Given $\mathcal{G}$ and $p,q$, compute an SLP for the string $\val(\mathcal{G})[p:q]$ (in case $0 \le p \leq q < |\val(\mathcal{G})|$ does not hold, the algorithm
outputs a special symbol).
\end{itemize}
\end{authorclaim}

\begin{authorclaim}{Lemma 33 (c.f.\ [38, Lemma 3.12])}  \label{lemma-SLP-iterated-morph}
Given a symbols $a_0 \in \Sigma$ and a sequence of morphisms $\phi_1, \ldots, \phi_n : \Sigma^* \to \Sigma^*$, where every $\phi_i$ is given by a list of the words
$\phi_i(a)$ for $a \in \Sigma$, one can compute in \L an SLP for the word
$\phi_1(\phi_2( \cdots \phi_n(a_0) \cdots))$.
\end{authorclaim}
The {\em circuit value problem} for a finitely generated group $G$ with the finite standard generating set $\Sigma$, $\CVP(G,\Sigma)$ for short,
is the following decision problem:
\begin{description}
\item[Input:]  an SLP $\mathcal{G}$ over the alphabet $\Sigma$.
\item [Question:] does $\val(\mathcal{G})=1$ hold in $G$?
\end{description}
It is an easy observation that the computational complexity of the compressed word 
problem for $G$ does not depend on the chosen generating set $\Sigma$ in the sense
that if $\Sigma'$ is another finite standard generating set for $G$, then $\CVP(G,\Sigma)$
is \L-reducible to $\CVP(G,\Sigma')$ \cite[Lemma~4.2]{ref38}. Therefore we do not
have to specify the generating set and we just write $\CVP(G)$.

The circuit value problem for $G$ is equivalent to the problem whether a given circuit over the group $G$ evaluates to $1$: Take an SLP $\mathcal{G} = (V,\rho,S)$ in
Chomsky normal form and built a circuit by taking $V$ is the set of gates. If $\rho(A)=a  \in \Sigma$ then $A$ is an input gate that is labelled
with the group generator $a$. If $\rho(A) = BC$ with $B,C \in V$ then $B$ is left input gate for $A$ and $C$ is the right input gate for $A$.
Such a circuit can be evaluated in the natural way (every internal gate computes the product of its input values) and the circuit output is the value
at gate $S$. 



\section*{Complete Appendix C proof}\setcounter{equation}{2}
This section presents a detailed proof of Theorem 6. Instead of circuits over groups we will use the equivalent framework of SLPs from Appendix B. The proof of the lower bound in Theorem 6 uses some of the techniques from the paper [37], where a connection between leaf strings and SLPs was established. In Sections C.1--C.3 we will introduce these techniques. The proof of Theorem 6 will be given in Appendix C.4.
Let us first fix some general notations. For integers $i\le j$ we write $[i..j]$ for the interval $\{z\in\mathbb Z\mid i\le z\le j\}$. In the following, we will identify a bit string $\alpha=a_1a_2\cdots a_n$ ($a_1,\ldots,a_n\in\{0,1\}$) with the vector $(a_1,a_2,\ldots,a_n)$. In particular, for another vector $\overline{s}=(s_1,\ldots,s_n)\in\mathbb N^n$ we will write $\alpha\cdot\overline{s}=\sum_{i=1}^n a_i\cdot s_i$ for the scalar product. Moreover, we write $\sum\overline{s}$ for the sum $s_1+\cdots+s_n$.
\subsection*{C.1 Subsetsum problems}  \label{sec-subsetsum}

A sequence $(s_1,\ldots,s_n)$ of natural numbers is
\emph{super-decreasing} if $s_i > s_{i+1}+ \cdots + s_n$ for all $i \in \ccinterval{1}{n}$. For example, $(s_1,\ldots,s_n)$
with $s_i=2^{n-i}$ is super-decreasing.
An instance of the \emph{subsetsum problem} is 
a tuple $(t,s_1,\ldots,s_k)$ of binary coded natural numbers.
It is a positive instance if there are $a_1, \ldots, a_k \in \{0,1\}$ such that
$t = a_1 s_1 + \cdots + a_k s_k$. 
Subsetsum is a classical $\NP$-complete problem, see e.g. \cite{ref15}.
The \emph{super-decreasing subsetsum} problem is the restriction 
of subsetsum to instances $(t,s_1,\ldots,s_k)$, where $(s_1, \ldots, s_k)$
is super-decreasing. In \cite{ref30} it was shown that super-decreasing 
subsetsum is $\Ptime$-complete.\footnote{In fact, \cite{ref30} deals with
the {\em super-increasing} subsetsum problem. But this is only a nonessential
detail. For our purpose, super-decreasing sequences are more convenient.} 
We need a slightly generalized version of the construction showing $\Ptime$-hardness 
that we discuss in Appendix C.2.


\subsection*{C.2 From Boolean circuits to super-decreasing subsetsum} \label{sec-circuits-subsetsum} 

For this section, we have to fix some more details on boolean circuits.
Let us consider a boolean circuit $C$ with input gates $x_1, \ldots, x_m$ 
and output gates $y_0, \ldots, y_{n-1}$.\footnote{It will be convenient for us to number the input gates from $1$ 
and the output gates from $0$.}
We view $C$ as a directed acyclic graph with multi-edges
(there can be two edges between two nodes); the nodes are the gates of the circuit.
The number of incoming edges of a gate is called its {\em fan-in} and the 
number of outgoing edges is the {\em fan-out}.
Every input gate $x_i$ has fan-in zero and every output gate $y_i$ has fan-out zero.
Besides the input gates there are two more gates $c_0$ and $c_1$ of fan-in zero, where
$c_i$ carries the constant truth value $i \in \{0,1\}$. Besides $x_1, \ldots, x_m, c_0, c_1$
every other gate has fan-in two and computes the \nand\ of its two input gates.
Moreover, we assume that every output gate $y_i$ is a {\sf nand}-gate.
For a bit string $\alpha = b_1 \cdots b_m$ ($b_1, \ldots, b_m \in \{0,1\}$) and $0 \leq i \leq n-1$
we denote with $C(\alpha)_i$ the value of the output gate $y_i$  when every input gate
$x_j$ ($1 \leq j \leq m$) is set to $b_j$. Thus, $C$ defines a map $\{0,1\}^m\to\{0,1\}^n$. 

We assume now that $C$ is a boolean circuit as above with the 
 following additional property that will be satisfied later: For all input bit strings $\alpha \in \{0,1\}^m$ there is exactly
one $i \in \cointerval{0}{n}$ such that $C(\alpha)_i=1$. Using a refinement of the construction from \cite{ref30}
we compute in \L $q_{0}, \ldots, q_{n-1} \in \mathbb{N}$ and two super-decreasing sequences 
$\overline{r} = (r_1, \ldots r_m)$ and $\overline{s} = (s_1, \ldots, s_k)$ for some $k$ (all numbers are represented in binary notation) with the following properties:
\begin{itemize}
\item The $r_1, \ldots, r_m$ are pairwise distinct powers of $4$. 
\item For all $0 \leq i \leq n-1$ and all $\alpha \in \{0,1\}^m$:
$C(\alpha)_i = 1$  if and only if there exists
$\delta \in \{0,1\}^k$ such that $\delta \cdot \overline{s} = q_i + \alpha \cdot \overline{r}$. 
\end{itemize}
Let us first add for every input gate $x_i$ two new {\sf nand}-gates
${\bar x}_i$ and ${\bar{\bar x}}_i$, where ${\bar{\bar x}}_i$ has the same outgoing edges as $x_i$. Moreover we remove the old outgoing edges of $x_i$
and replace them by the edges $(x_i, {\bar x}_i)$, $(c_1, {\bar x}_i)$ and two edges from ${\bar x}_i$ to ${\bar{\bar x}}_i$.
This has the effect that every input gate $x_i$ has a unique outgoing edge. Clearly, the new circuit computes the same
boolean function (basically, we introduce two negation gates for every input gate).
Let $g_1, \ldots, g_p$ be the {\sf nand}-gates of the circuit enumerated in reverse topological
order, i.e., if there is an edge from gate $g_i$ to gate $g_j$ then $i>j$. We denote the two edges entering gate
$g_i$ with $e_{2i+n-2}$ and $e_{2i+n-1}$. Moreover, we write $e_i$ ($0 \leq i \leq n-1$) for an imaginary edge
that leaves the output gate $y_{i}$ and whose target gate is unspecified. Thus, the edges of the circuit are 
$e_0,\ldots, e_{2p+n-1}$.
We now define the natural numbers $q_{0}, \ldots, q_{n-1}, r_1, \ldots r_m,s_1, \ldots, s_k$ with $k = 3p$:
\begin{itemize}
\item
Let $I = \{ j \mid e_j \text{ is an outgoing edge of the constant gate  $c_1$ or a {\sf nand}-gate} \}$.
For $0 \leq i \leq n-1$ we define the number $q_{i}$ as 
$$
q_{i} = \sum_{j \in I \setminus \{i\}} 4^j .
$$
Recall that $e_i$ is the unique outgoing edge of the output gate $y_i$.
\item If $e_j$ is the unique outgoing edge of the input gate $x_i$ then we set $r_i = 4^j$. We can choose the reverse
topological sorting of the {\sf nand}-gates in such a way that $r_1 > r_2 > \cdots > r_m$ (we only have to ensure that the target
gates $\overline{x}_1, \ldots, \overline{x}_m$ of the input gates appear in the order $\overline{x}_m, \ldots, \overline{x}_1$ in the 
reverse topological sorting of the {\sf nand}-gates).
\item To define the numbers $s_1, \ldots, s_k$ we first define for every {\sf nand}-gate $g_i$
three numbers $t_{3i}$, $t_{3i-1}$ and $t_{3i-2}$ as follows, where 
$I_i = \{ j \mid e_j \text{ is an outgoing edge of gate } g_i \}$:
\begin{eqnarray*}
t_{3i}       &=& 4^{2i+n-1} + 4^{2i+n-2} + \sum_{j \in I_i} 4^j \\
t_{3i-1} &=& 4^{2i+n-1} - 4^{2i+n-2} = 3 \cdot 4^{2i+n-2} \\
t_{3i-2}    &=& 4^{2i+n-2}   
\end{eqnarray*}
Then, the tuple $(s_1,\ldots, s_k)$ is $(t_{3p}, t_{3p-1}, t_{3p-2}, \ldots, t_3, t_2, t_1)$, which is indeed
super-decreasing (see also \cite{ref30}). In fact, we have $s_i - (s_{i+1} + \cdots + s_k) \geq 4^{n-1}$ for all $i \in \ccinterval{1}{k}$.
To see this, note that the sets $I_{i+1},\ldots, I_k$ are pairwise disjoint.  This implies that the $n-1$ low-order digits in the base-4 expansion of 
$s_{i+1} + \cdots + s_k$ are zero or one.

\end{itemize}
In order to understand this construction, one should think of the edges of the circuit carrying truth values.
Recall that there are $2p+n$ edges in the circuit (including the imaginary outgoing edges of the output gates
$y_0, \ldots, y_{n-1}$). 
A number in base-4 representation with $2p+n$ digits that are either $0$ or $1$ represents a truth assignment to the $2p+n$ edges,
where a $1$-digit represents the truth value $1$ and a $0$-digit represents the truth value $0$.
Consider  an input string $\alpha = b_1 \cdots b_m \in \{0,1\}^m$  and 
consider an output gate $y_i$, $i \in \cointerval{0}{n}$. Then the number
$N := 4^i + q_{i} + b_1 r_1 + \cdots + b_m r_m$ encodes the truth assignment for the circuit edges, where:
\begin{itemize}
\item all outgoing edges of the constant gate $c_1$ carry the truth value $1$, 
\item all outgoing edges of the constant gate $c_0$ carry the truth value $0$, 
\item the unique outgoing edge of an input gate $x_i$ carries the truth value $b_i$,
\item all outgoing edges of {\sf nand}-gates 
carry the truth value $1$.
\end{itemize}
We have to show that $C(\alpha)_i = 1$  if and only if there exists
$\delta \in \{0,1\}^k$ such that $\delta \cdot \overline{s} = N-4^i$.
For this we apply the canonical algorithm for super-decreasing subsetsum with input $(N,s_1,\ldots,s_k)$.
This algorithm initializes a counter $A$ to $N$ and then goes over the sequence $s_1,\ldots,s_k$ in that order.
In the $j$-th step ($1 \leq j \leq k$) we set $A$ to $A-s_j$ if $A \geq s_j$. If $A < s_j$ then we do not modify $A$.
After that we proceed with $s_{j+1}$. The point is that 
this process simulates the evaluation of the circuit on the input values $b_1, \ldots, b_m$. Thereby
the {\sf nand}-gates are evaluated in the topological order $g_p, g_{p-1}, \ldots, g_1$. Assume that $g_j$ is the gate that we 
want to evaluate next. In the above algorithm for super-decreasing subsetsum the evaluation of $g_j$ is simulated 
by the three numbers $t_{3j}$, $t_{3j-1}$, and $t_{3j-2}$. At the point where the algorithm checks whether $t_{3j}$ can 
be subtracted from $A$, the base-4 digits at positions $2j+n,\ldots, 2p+n-1$ in the counter value $A$ have been already set to zero. If the digits at the next two high-order 
positions $2j+n-1$ and $2j+n-2$ are still $1$ (i.e., the input edges $e_{2j+n-2}$ and $e_{2j+n-1}$ for gate
$g_j$ carry the truth value $1$), then we can subtract $t_{3j}$ from $A$. Thereby we subtract all powers
$4^{2j+n-1}$, $4^{2j+n-2}$ and
$4^h$, where $e_h$ is an outgoing edge for gate $g_j$. Since gate $g_j$ evaluates to zero (both input edges carry $1$), this subtraction
correctly simulates the evaluation of gate $g_j$: all outgoing edges $e_h$ of $g_j$ (that were initially set to the truth value $1$) are set to the 
truth value $0$. On the other hand, if one of the two digits at
positions $2j+n-1$ and $2j+n-2$ in $A$ is $0$ (which means that gate $g_j$ evaluates to $1$), then we 
cannot subtract $t_{3j}$ from $A$. If both digits at positions $2j+n-1$ and $2j+n-2$ in $A$ are $0$, then
also $t_{3j-1}$ and $t_{3j-2}$ cannot be subtracted. On the other hand, if exactly one of the two digits
at positions $2j+n-1$ and $2j+n-2$ is $1$, then with
$t_{3j-1}$ and $t_{3j-2}$ we can set  these two digits to $0$ (thereby digits at positions $< 2j+n-2$ are not modified).

Assume now that $y_j$ ($j \in \cointerval{0}{n}$) is the unique output gate that evaluates to $1$, i.e., all
output gates $y_{j'}$ with $j' \neq j$ evaluate to zero. Then after processing all weights
$s_1,\ldots,s_k$ we have $A = 4^j$ (we will never subtract $4^j$). We have shown that there exists
$\delta \in \{0,1\}^k$ such that $\delta \cdot \overline{s} + 4^j = N$. Hence, if $i=j$ (i.e., $C(\alpha)_i = 1$)
then  $\delta \cdot \overline{s} = N-4^i$. Now assume that $i \neq j$. In order to get a contradiction assume 
that there is $\delta' \in \{0,1\}^k$ such that $\delta' \cdot \overline{s} + 4^i = N$. We have $\delta \neq \delta'$ and 
$\delta \cdot \overline{s} + 4^j  = \delta' \cdot \overline{s} + 4^i$, i.e, $\delta \cdot \overline{s} - \delta' \cdot \overline{s} = 4^i - 4^j$.
Since $i,j \in \cointerval{0}{n}$ we get $|\delta \cdot \overline{s} - \delta' \cdot \overline{s}| < 4^{n-1}$. But 
$s_i - (s_{i+1} + \cdots s_k) \geq 4^{n-1}$ for all $i \in \ccinterval{1}{k}$ implies that 
$|\delta \cdot \overline{s} - \delta' \cdot \overline{s}| \ge 4^{n-1}$.

\subsection*{C.3 From super-decreasing subsetsum to straight-line programs} \label{sec-subsetsum-SLP}

In \cite{ref35} a super-decreasing sequence $\overline{t} = (t_1,\ldots,t_k)$ of natural numbers is encoded by the string 
$S(\overline{t}) \in \{0,1\}^*$ of length $\sum \overline{t} +1$
such that for all $0 \leq p \leq \sum \overline{t}$:
\begin{equation}\label{main string}
S(\overline{t})[p] = \begin{cases} 1 & \text{if $p = \alpha \cdot \overline{t}$ for some $\alpha  \in \{0,1\}^k$}, \\
0 & \text{otherwise}.
\end{cases}
\end{equation}
Note that in the first case, $\alpha$ is unique.
Since $\overline{t}$ is a super-decreasing sequence, 
the number of $1$'s in the string $S(\overline{t})$ is $2^k$.
Also note that $S(\overline{t})$ starts and ends with $1$.
In \cite{ref35} it was shown that from a super-decreasing sequence $\overline{t}$
of binary encoded numbers one can construct in \L an SLP for 
the word $S(\overline{t})$.



   
\subsection*{C.4 Proof of Theorem 6} \label{sec:leaf-wreath}

Let us fix a regular wreath product of the form $G \wr \mathbb{Z}$ for a finitely generated
group $G$. Such groups are also known as generalized lamplighter groups (the lamplighter group arises
for $G = \mathbb{Z}_2$). Throughout this section, we fix a set of standard generators $\Sigma$ for $G$ and let $\tau = 1$ be the generator for $\mathbb{Z}$. 
Then $\Sigma \cup \{\tau,\tau^{-1}\}$ is a standard generating set for the wreath product $G \wr \mathbb{Z}$.
In $G \wr \mathbb{Z}$ the $G$-generator $a \in \Sigma$ represents the mapping $f_{a} \in G^{(\Z)}$ with $f_{a}(0) = a$ and 
$f_{a}(z) = 1$ for $z \neq 0$. 
For a word $w \in (\Sigma \cup \{\tau,\tau^{-1}\})^*$ we define $\eta(w) :=  |w|_\tau-|w|_{\tau^{-1}}$. Thus, the element of $G \wr \mathbb{Z}$ represented by
$w$ is of the form $f  \tau^{\eta(w)}$ for some $f \in G^{(\Z)}$.
Recall the definition of the left action of $\Z$ on $G^{(\Z)}$ from the wreath-product conventions (where we take $H = Y = \Z$).
For better readability, we write $c \circ f$ for ${}^c \! f$ ($c \in \Z$, $f \in G^{(\Z)}$).
Hence, we have $(c \circ f)(z) = f(z+c)$. If one thinks of $f$ as a bi-infinite word over the alphabet $G$,
then $c \circ f$ is the same word but shifted by $-c$.

The following intuition might be helpful: Consider a word $w \in ( \Sigma \cup \{\tau, \tau^{-1} \})^*$. In 
$G \wr \mathbb{Z}$ we can simplify $w$ to a word of the form
$\tau^{z_0} a_{1} \tau^{z_1} a_{2} \cdots \tau^{z_{k-1}} a_{k} \tau^{z_k}$ (with $z_j \in \Z$, $a_j \in \Sigma$),
which in $G \wr \mathbb{Z}$ can be rewritten
as
\[
\tau^{z_0} a_{1} \tau^{z_1} a_{2} \cdots \tau^{z_{k-1}} a_{k} \tau^{z_k} = \big(\prod_{j=1}^k \tau^{z_0 + \cdots+ z_{j-1}} a_{j} \tau^{-(z_0 + \cdots +z_{j-1})}\big) \; \tau^{z_0 + \cdots +z_k}.
\]  
Hence, the word $w$ represents the group element 
\[
\big(\prod_{j=1}^k (z_0 + \cdots+ z_{j-1}) \circ f_{a_j} \big)\, \tau^{z_0 + \cdots +z_k} .
\]
This gives the following intuition for evaluating $\tau^{z_0} a_{1} \tau^{z_1} a_{2} \cdots \tau^{z_{k-1}} a_{k} \tau^{z_k}$:
In the beginning, every $\Z$-position carries the $G$-value $1$. First, 
go to the $\Z$-position $-z_0$ and multiply the $G$-element at this position with $a_1$ (on the right), then 
go to the $\Z$-position $-z_0-z_1$ and multiply the $G$-element at this position with $a_2$, and so on.  

\begin{proof}[Proof of Theorem 6]
The easy part is to show that the circuit value problem for $G \wr \mathbb{Z}$ belongs to
$\forall\LEAF(\WP(G))$. In the following, we make use
of the statements from Lemma 32.
Let $\mathcal{G}$ be an SLP over the alphabet $\Sigma \cup \{\tau,\tau^{-1}\}$ and let $f \tau^{\eta(\val(\mathcal{G}))} \in G \wr \mathbb{Z}$
be the group element represented by $\val(\mathcal{G})$. By Lemma 32 we can compute $\eta(\val(\mathcal{G}))$ in polynomial time. 
If $\eta(\val(\mathcal{G})) \neq 0$ then the Turing-machine rejects by printing a non-trivial generator of $G$ (here we need the assumption that
$G$ is non-trivial).
So, let us assume that $\eta(\val(\mathcal{G}))=0$. 
We can also compute in polynomial time two integers $b,c \in \mathbb{Z}$ such that $\supp(f) \subseteq \ccinterval{b}{c}$. We can
take for instance $b = -|\val(\mathcal{G})|$ and $c = |\val(\mathcal{G})|$. It suffices to check whether for all
$ x \in \ccinterval{b}{c}$ we have $ f(x)=1$. For this, the Turing-machine branches universally to all binary coded integers $x \in \ccinterval{b}{c}$
(this yields the $\forall$-part in $\forall \LEAF(\WP(G))$). Consider a specific branch that leads to the integer 
$x \in \ccinterval{b}{c}$. From $x$ and the input SLP $\mathcal{G}$ the Turing-machine then produces a leaf string over the standard
generating set $\Sigma$ of $G$ such that this leaf string represents the group element $f(x) \in G$. For this, the machine
branches to all positions $p \in \cointerval{0}{|\val(\mathcal{G})|}$ (if $p<q<|\val(\mathcal{G})|$ then the branch for $p$ is to the left 
of the branch for $q$). For a specific position $p$, the machine computes in polynomial time the symbol
$a = \val(\mathcal{G})[p]$. If $a$ is $\tau$ or $\tau^{-1}$ then the machine prints $1 \in \Sigma$. On the other hand, if 
$a \in \Sigma$ then the machine computes in polynomial time $d = \eta(\val(\mathcal{G})[:p])$. This is possible by first
computing an SLP for the prefix $\val(\mathcal{G})[:p]$. If $d=-x$ then the machine prints the symbol $a$, otherwise
the machine prints the trivial generator $1$. It is easy to observe that the leaf string produced in this way represents the 
group element $f(x)$.

We now show the hardness statement from Theorem 6.
By Lemma 29 it suffices to show that $\CVP(G \wr \mathbb{Z})$ 
is hard for $\forall\bLEAF(\WP(G/Z(G)))$ with respect to \L-reductions. Let $a_0, \ldots, a_{n-1}$ be an arbitrary enumeration of the standard generators in $\Sigma$.
Fix a language $L \in \forall\bLEAF(\WP(G/Z(G)))$. From the definition of the class $\forall\bLEAF(\WP(G/Z(G)))$ it follows that
there exist two polynomials $p_1$ and $p_2$ and a balanced polynomial time NTM $M$ running in time $p_1+p_2$
that outputs a symbol from $\Sigma$ after termination and such that the following holds: Consider an input word
$z$ and let $T(z)$ be the corresponding computation tree of $M$. 
Let $m_1 = p_1(|z|)$, $m_2 = p_2(|z|)$, and $m = m_1+m_2$. Note that the nodes of $T(z)$ are the bit strings of length at most $m$.
For every leaf $\alpha \in \{0,1\}^m$ let us denote with $\lambda(\alpha)$ the symbol from $\Sigma$ that $M$ prints 
when reaching the leaf $\alpha$. Then $z \in L$ if and only if for all
$\beta \in \{0,1\}^{m_1}$ the string 
\begin{equation} \label{eq-alpha-u}
\lambda_\beta := \prod_{\gamma \in \{0,1\}^{m_2}}\lambda(\beta\gamma)
\end{equation}
represents a group element from the center $Z(G)$. Here (and in the following), the product in the right-hand side of  \eqref{eq-alpha-u} goes over all bit strings of length $m_2$
in lexicographic order.
Our construction consists of five steps:

\medskip
\noindent
{\em Step 1.}
Note that given a bit string $\alpha \in \{0,1\}^m$, we can 
compute in polynomial time the symbol $\lambda(\alpha) \in \Sigma$ by following the computation path specified by $\alpha$. 
Using the classical Cook-Levin construction 
(see e.g. \cite{ref2}), we can compute from the input $z$ and $a \in \Sigma$  in \L
a boolean circuit $C_{z,a}$ with $m$ input gates $x_1, \ldots, x_m$ and a single output gate $y_0$
such that for all $\alpha \in \{0,1\}^m$:
$C_{z,a}(\alpha)_0 = 1$ if and only if $\lambda(\alpha) = a$. 
By taking the disjoint union of these circuits and merging the input gates,
we can build a single circuit $C_z$ with $m$ input gates $x_1, \ldots, x_m$
and $n = |\Sigma|$ output gates $y_0, \ldots, y_{n-1}$. For every $\alpha \in \{0,1\}^m$ and every
$0 \leq i \leq n-1$ the following holds: $C_z(\alpha)_i = 1$ if and only if $\lambda(\alpha) = a_i$.

\medskip
\noindent
{\em Step 2.}
Using the construction from Appendix C.2 we can compute from the  
circuit $C_z$ in \L numbers $q_{0}, \ldots, q_{n-1} \in \mathbb{N}$ and two super-decreasing sequences 
$\overline{r} = (r_1, \ldots, r_m)$ and $\overline{s} =(s_1, \ldots, s_k)$ with the following properties:
\begin{itemize}
\item The $r_1, \ldots, r_m$ are pairwise distinct powers of $4$. 
\item For all $0 \leq i \leq n-1$ and all $\alpha \in \{0,1\}^m$ we have:
$\lambda(\alpha) = a_i$ if and only if
$C_z(\alpha)_i = 1$  if and only if there is $\delta \in \{0,1\}^k$ such that $\delta \cdot \overline{s} = q_{i} + \alpha \cdot \overline{r}$. 
\end{itemize}
Note that for all $\alpha \in \{0,1\}^m$ there is a unique $i$ such that $C_z(\alpha)_i = 1$.
Hence, for all  $\alpha \in \{0,1\}^m$ there is a unique $i$ such that $q_{i} + \alpha \cdot \overline{r}$ is of the form 
$\delta \cdot \overline{s}$ for some $\delta \in \{0,1\}^k$. For this unique $i$ we have $\lambda(\alpha) = a_i$.

We split the super-decreasing sequence $\overline{r} = (r_1, \ldots, r_m)$ into the two
sequences $\overline{r}_1 = (r_1, \ldots, r_{m_1})$ and $\overline{r}_2 = (r_{m_1+1}, \ldots, r_m)$.
For the following consideration, we need the following numbers:
\begin{eqnarray}
\ell & = & \textstyle \max\big\{ \sum \overline{r}_1 + \max\{q_0, \ldots, q_{n-1}\}+1, \; \sum\overline{s} - \sum\overline{r}_2 - \min \{q_0, \ldots, q_{n-1}\}+1\big\} \label{eq:l} \\
\pi &= & \textstyle \ell + \sum \overline{r}_2 
\end{eqnarray}
The binary codings of these numbers can be computed in \L (since iterated addition, $\max$, and $\min$ can be computed in \L).
The precise value of $\ell$ will be only relevant at the end of step 4.

\medskip
\noindent
{\em Step 3.} 
By the result from \cite{ref35} (see Section~\ref{sec-subsetsum-SLP})
we can construct in \L from the three super-decreasing sequences
$\overline{r}_1, \overline{r}_2$ and $\overline{s}$ three SLPs $\mathcal{G}_1$,   $\mathcal{G}_2$ and $\mathcal{H}$ over the alphabet
$\{0,1\}$ such that $\val(\mathcal{G}_1) = S(\overline{r}_1)$, $\val(\mathcal{G}_2) = S(\overline{r}_2)$
and $\val(\mathcal{H}) = S(\overline{s})$ (see (\ref{main string})).
For all positions $p \geq 0$ (in the suitable range) we have:
\begin{eqnarray*}
\val(\mathcal{G}_1)[p] = 1  & \Longleftrightarrow & 
\exists \beta \in  \{0,1\}^{m_1}: 
p = \beta \cdot \overline{r}_1 \\ 
\val(\mathcal{G}_2)[p] = 1  & \Longleftrightarrow & 
\exists  \gamma \in  \{0,1\}^{m_2}: 
p = \gamma \cdot \overline{r}_2 \\ 
\val(\mathcal{H})[p] = 1  & \Longleftrightarrow & 
\exists \delta \in  \{0,1\}^k: 
p = \delta \cdot \overline{s}
\end{eqnarray*}
Note that $|\val(\mathcal{G}_1)| = \sum\overline{r}_1+1$, $|\val(\mathcal{G}_2)| = \sum\overline{r}_2+1$, and
$|\val(\mathcal{H})| = \sum\overline{s}+1$.


\medskip
\noindent
{\em Step 4.}
We build in \L for every $i \in \ccinterval{0}{n-1}$ an
SLP $\mathcal{H}_i$ 
from the SLP $\mathcal{H}$ by replacing in every right-hand side of $\mathcal{H}$ 
every occurrence of $0$  by $\tau^{-1}$ and
every occurrence of $1$  by $a_i \tau^{-1}$.
Let $T_i$ be the start variable of $\mathcal{H}_i$, let $S_1$ be the start variable of $\mathcal{G}_1$, and let
$S_2$ be the start variable of $\mathcal{G}_2$.
We can assume that the variable sets of the SLPs $\mathcal{G}_1, \mathcal{G}_2, \mathcal{H}_0, \ldots, \mathcal{H}_{n-1}$
are pairwise disjoint.  We next combine these SLPs into a single SLP $\mathcal{I}$.
The variables of $\mathcal{I}$ are the variables of 
the SLPs $\mathcal{G}_1, \mathcal{G}_2,  \mathcal{H}_0, \ldots, \mathcal{H}_{n-1}$ plus a fresh variable $S$ which is the start
variable of $\mathcal{I}$. The right-hand sides for the variables are defined below. In the right-hand sides 
we write powers $\tau^p$ for integers $p$
whose binary codings can be computed in \L. Such powers can be produced by small subSLPs that can be constructed in \L too.
\begin{itemize}
\item In all right-hand sides of $\mathcal{G}_1$  and $\mathcal{G}_2$ 
we replace all occurrences of the terminal symbol $0$ by the $\mathbb{Z}$-generator $\tau$. 
\item We replace every
occurrence of the terminal symbol $1$ in a right-hand side of $\mathcal{G}_1$ by $S_2 \tau^\ell$,
where $\ell$ is from \eqref{eq:l}.
\item We  replace every
occurrence of the terminal symbol $1$ in a right-hand side of $\mathcal{G}_2$ by
$\sigma \tau$, where
\begin{equation} \label{eq-sigma}
\sigma = \tau^{q_{0}} T_0 \tau^{h-q_{0}} \tau^{q_{1}} T_1 \tau^{h-q_{1}} \cdots \tau^{q_{n-1}} T_{n-1} \tau^{h-q_{n-1}}
\end{equation}
and $h = \sum \overline{s}+1$ is the length of the word
$\val(\mathcal{H})$ (which is $-\eta(\val_{\mathcal{I}}(T_i))$ for every  $i \in\cointerval{0}{n}$).
Note that $\eta(\val_{\mathcal{I}}(\sigma)) = 0$.
\item Finally, 
the right-hand side of the start variable $S$ is  $S_1 \tau^{-d}$ where $d := \sum \overline{r}_1 + 1 + 2^{m_1} \cdot \pi$.
(note that $d = \eta(\val_{\mathcal{I}}(S_1))$).
\end{itemize}
Before we explain this construction, let us first introduce some notations.
\begin{itemize}
\item Let $u := \val_{\mathcal{I}}(S_2)$. We have $\eta(u) = |\val(\mathcal{G}_2)|$. Hence,
the group element represented by $u$ can be written as
$f_u \tau^{|\val(\mathcal{G}_2)|}$ for a mapping $f_u \in G^{(\Z)}$. 
\item Let $v := \val_{\mathcal{I}}(\sigma)$ where $\sigma$ is from \eqref{eq-sigma}. Note that $\eta(v)=0$.
Hence, the group element represented by $v$ is a mapping $f_v \in G^{(\Z)}$. Its support is a subset of 
the interval from position $-\max\{ q_0, \ldots, q_{n-1}\}$ 
to position $\sum \overline{s}-\min\{ q_0, \ldots, q_{n-1}\}$.
\item For $\beta \in \{0,1\}^{m_1}$ let
$\mathsf{bin}(\beta)$ be the number represented by $\beta$ in binary notation (thus, $\mathsf{bin}(0^{m_1}) = 0$, $\mathsf{bin}(0^{m_1-1}1) = 1$, \ldots, 
$\mathsf{bin}(1^{m_1}) = 2^{m_1}-1$). Moreover, let 
\[ p_\beta := -\mathsf{bin}(\beta) \cdot \pi.\]
\end{itemize}
First, note that $\eta(\val(\mathcal{I})) = 0$. This is due to the factor $\tau^{-d}$ in the right-hand side of the start
variable $S$ of $\mathcal{I}$. Hence, the group element represented by $\val(\mathcal{I})$ is a mapping $f \in G^{(\Z)}$.
The crucial claim is the following:

\medskip
\noindent
\begin{authorclaim}{Claim}
	For every $\beta \in \{0,1\}^{m_1}$, $f(p_{\beta})$ is the group element represented by the leaf string
	$\lambda_{\beta}$ from \eqref{eq-alpha-u}.
\end{authorclaim}


\medskip
\noindent
{\em Proof of the claim.} In the following, we compute in the restricted direct product $G^{(\mathbb{Z})}$.
Recall that the multiplication in this group is defined by the pointwise multiplication of mappings.

Since we replaced in $\mathcal{G}_1$ every $1$ in a right-hand side by 
$S_2 \tau^\ell$, which produces $u \tau^\ell$ in $\mathcal{I}$ (which evaluates to $f_u \tau^{\pi+1}$)
the mapping $f$ is a product (in the restricted direct product $G^{(\mathbb{Z})}$) of shifted copies of $f_u$.
More precisely, for every $\beta' \in \{0,1\}^{m_1}$ we get the shifted copy 
\begin{equation} \label{eq-shift-dist}
 \big(\beta' \cdot \overline{r}_1 + \mathsf{bin}(\beta') \cdot \pi \big) \circ f_u 
\end{equation}
of $f_u$. The shift distance $\beta' \cdot \overline{r}_1 + \mathsf{bin}(\beta') \cdot \pi$ 
can be explained as follows: The $1$ in $\val(\mathcal{G}_1)$ that corresponds to $\beta' \in \{0,1\}^{m_1}$
occurs at position $\beta' \cdot \overline{r}_1$ (the first position is $0$) and to the left of this position we find $\mathsf{bin}(\beta')$ many
$1$'s and $\beta' \cdot \overline{r}_1 - \mathsf{bin}(\beta')$ many $0$'s in $\val(\mathcal{G}_1)$.
Moreover, every $0$ in $\val(\mathcal{G}_1)$ was replaced
by $\tau$ (shift by $1$) and every $1$ in $\val(\mathcal{G}_1)$ was replaced by $u \tau^\ell$ 
(shift by $\ell + |\val(\mathcal{G}_2)| = \pi + 1$). Hence, the total shift distance
is indeed  \eqref{eq-shift-dist}. Also note that
if $\beta' \in \{0,1\}^{m_1}$ is lexicographically smaller than $\beta'' \in \{0,1\}^{m_1}$ then 
$\beta' \cdot \overline{r}_1 < \beta'' \cdot \overline{r}_1$. 
This implies that 
\begin{eqnarray*}
f = \prod_{\beta' \in \{0,1\}^{m_1}}  \big(\beta' \cdot \overline{r}_1 + \mathsf{bin}(\beta') \cdot \pi \big) \circ f_u  = 
\prod_{\beta' \in \{0,1\}^{m_1}}  \big(\beta' \cdot \overline{r}_1 -p_{\beta'}  \big) \circ f_u .
\end{eqnarray*}
Let us now compute the mapping $f_u$. Recall that we replaced in $\mathcal{G}_2$ every occurrence of $1$ by  $\sigma \tau$, where $\sigma$ is from
\eqref{eq-sigma} and derives to $v$. The $1$'s in $\val(\mathcal{G}_2)$ occur at positions of the form $\gamma \cdot \overline{r}_2$
for $\gamma \in \{0,1\}^{m_2}$ and if $\gamma \in \{0,1\}^{m_2}$ is lexicographically smaller than $\gamma' \in \{0,1\}^{m_2}$ then 
$\gamma \cdot \overline{r}_2 < \gamma' \cdot \overline{r}_2$.
We therefore get
\[
f_u = \prod_{\gamma \in \{0,1\}^{m_2}}  (\gamma \cdot \overline{r}_2) \circ f_v . 
\]
We obtain 
\begin{eqnarray*}
f &=&  \prod_{\beta' \in \{0,1\}^{m_1}}  \big(\beta' \cdot \overline{r}_1 - p_{\beta'}\big) \circ f_u  \\
& = & \prod_{\beta' \in \{0,1\}^{m_1}}   \big(\beta' \cdot \overline{r}_1 - p_{\beta'} \big) \circ \prod_{\gamma \in \{0,1\}^{m_2}} ( \gamma \cdot \overline{r}_2 \circ f_v)   \\
& = & \prod_{\beta' \in \{0,1\}^{m_1}}  \prod_{\gamma \in \{0,1\}^{m_2}}   \big(\beta' \cdot \overline{r}_1 + \gamma \cdot \overline{r}_2 - p_{\beta'} \big) \circ f_v
\end{eqnarray*}
and hence
\begin{eqnarray*}
f(p_\beta) 
&=& \prod_{\beta' \in \{0,1\}^{m_1}}  \prod_{\gamma \in \{0,1\}^{m_2}} f_v(p_{\beta} - p_{\beta'} + \beta' \cdot \overline{r}_1 + \gamma \cdot \overline{r}_2) .
\end{eqnarray*}
We claim that for all $\beta \neq \beta'$ and all $\gamma\in \{0,1\}^{m_2}$ we have
\begin{equation} \label{eq-f_v}
f_v(p_{\beta} - p_{\beta'} + \beta' \cdot \overline{r}_1 + \gamma \cdot \overline{r}_2) = 1 .
\end{equation}
Let us postpone the proof of this for a moment. From \eqref{eq-f_v} we get
\[
f(p_\beta) = \prod_{\gamma \in \{0,1\}^{m_2}} f_v(\beta \cdot \overline{r}_1 + \gamma \cdot \overline{r}_2) .
\]
Consider a specific $\gamma \in \{0,1\}^{m_2}$ and let $\alpha = \beta \gamma$ and
$p = \beta \cdot \overline{r}_1 + \gamma \cdot \overline{r}_2 = \alpha \cdot \overline{r}$.
From the definition of $v = \val_{\mathcal{I}}(\sigma)$ it follows that for all $x \in \Z$, $f_v(x)$ is a product
of those group generators $a_i$ such that $x =  -q_i + \delta \cdot \overline{s}$ for some $\delta \in \{0,1\}^k$.
For the position $p$ this means that $q_i + \alpha \cdot \overline{r} = \delta \cdot \overline{s}$. By our previous remarks,
there is a unique such $i \in\cointerval{0}{n}$ and for this $i$ we have $\lambda(\alpha) = a_i$.
Hence, we obtain $f_v(p) = \lambda(\alpha)= \lambda(\beta\gamma)$ and thus
\[
f(p_\beta) = \prod_{\gamma \in \{0,1\}^{m_2}} \lambda(\beta\gamma) = \lambda_\beta .
\]
It remains to show \eqref{eq-f_v}. To get this identity, we need the precise value of $\ell$ from \eqref{eq:l}
(so far, the  value of $\ell$ was not relevant).
Assume now that $\beta \neq \beta'$, which implies
\[
|p_{\beta} - p_{\beta'}| \geq \pi =  \ell +  \sum \overline{r}_2 .
\]
Hence, we  either have 
\begin{eqnarray*}
p_{\beta} - p_{\beta'} +\beta' \cdot \overline{r}_1 + \gamma \cdot \overline{r}_2 & \geq & \ell +  \sum \overline{r}_2 + \beta' \cdot \overline{r}_1 + \gamma \cdot \overline{r}_2 \\
& \geq & \ell + \sum \overline{r}_2 \\
& > & \sum \overline{s} -  \min\{ q_0, \ldots, q_{n-1}\}
\end{eqnarray*}
or
\begin{eqnarray*}
p_{\beta} - p_{\beta'} +\beta' \cdot \overline{r}_1 + \gamma \cdot \overline{r}_2 & \leq & -\ell - \sum \overline{r}_2 + \beta' \cdot \overline{r}_1 + \gamma \cdot \overline{r}_2 \\
& \leq & -\ell + \sum \overline{r}_1  \\
& < &  - \max\{ q_0, \ldots, q_{n-1}\},
\end{eqnarray*}
where the strict inequalities follow from our choice of $\ell$.
Recall that the support of the mapping $f_v$ is contained  in $\ccinterval{-\max\{ q_0, \ldots, q_{n-1}\}}{\sum \overline{s}-\min\{ q_0, \ldots, q_{n-1}\}}$.
This shows \eqref{eq-f_v} and hence the claim.



\medskip
\noindent
{\em Step 5.} By the above claim, we have $f(p_{\beta}) \in Z(G)$ for all $\beta \in \{0,1\}^{m_1}$ if and only if 
$\lambda_{\beta} \in Z(G)$ for all $\beta \in \{0,1\}^{m_1}$, which is equivalent to $z \in L$. 
The only remaining problem is that the word $\val(\mathcal{I})$ produces some ``garbage'' group elements $f(x)$ on 
positions $x$ that are not of the form $p_\beta$. Note that for every $g \in G \setminus Z(G)$, there is a generator $a_i \in \Sigma$ such that the commutator 
$[g,a_i]$ is non-trivial. We now produce from $\mathcal{I}$ an SLP $\mathcal{I}^{-1}$ such that $\val(\mathcal{I}^{-1})$ represents
the inverse element of $f \in G^{(\Z)}$, which is the mapping $g$ with $g(x) = f(x)^{-1}$ for all $x \in \mathbb{Z}$. To construct $\mathcal{I}^{-1}$,
we have to reverse every right-hand side of $\mathcal{I}$ and replace every occurrence of a symbol $a_0, \ldots, a_{n-1}, \tau, \tau^{-1}$ by
its inverse. 

It is easy to compute in \L for every $i \in \cointerval{0}{n}$ an SLP for the word 
\[ 
w_i := \big(a_i \tau^\pi \big)^{2^{m_1}} \tau^{-2^{m_1} \cdot \pi} .
\]
Then the group element represented by $w_i$ is the mapping $f_i \in G^{(\Z)}$  whose  support 
is the set of positions $p_\beta$ for $\beta \in \{0,1\}^{m_1}$ and 
$f_i(p_\beta) = a_i$ for all $\beta \in \{0,1\}^{m_1}$.
We can also compute in \L an SLP for the word $w_i^{-1}$.
We then built in \L SLPs $\mathcal{J}_0, \ldots, \mathcal{J}_{n-1}$ such that $\val(\mathcal{J}_i) = \val(\mathcal{I}^{-1}) w_i^{-1} \val(\mathcal{I}) w_i$.
Hence, the word $\val(\mathcal{J}_i)$ represents the group element $g_i \in G^{(\Z)}$, where $g_i(x) = 1$ for all 
$x \in \mathbb{Z}\setminus \{ p_\beta \mid \beta \in \{0,1\}^{m_1} \}$ and 
$g_i(p_\beta) = f(p_\beta)^{-1} a_i^{-1} f(p_\beta) a_i  = [f(p_\beta),a_i]$. 

Finally, we construct in \L an SLP $\mathcal{J}$ such that 
\[ 
\val(\mathcal{J}) = \val(\mathcal{J}_0)\, \tau \,\val(\mathcal{J}_1) \,\tau\, \val(\mathcal{J}_2) \cdots \tau \, \val(\mathcal{J}_{n-1}) \,\tau^{-n+1} .
\]
We can assume that $n \leq \ell+\sum \overline{r}_2 = \pi$ ($n$ is a constant and we can always make $\ell$ bigger).
Then $\val(\mathcal{J})$ evaluates to the group element $g \in G^{(\Z)}$ with $g(x) = 1$ for 
$x \in \mathbb{Z} \setminus \{ p_\beta - i \mid \beta \in \{0,1\}^{m_1}, 0 \leq i \leq n-1 \}$ and
$g(p_\beta - i) = g_i(p_\beta) = [f(p_\beta),a_i]$ for $0 \leq i \leq n-1$.
Hence, if $f(p_{\beta}) \in Z(G)$ for all $\beta \in \{0,1\}^{m_1}$ then 
$\val(\mathcal{J})=1$ in $G \wr \Z$. On the other hand, if there is a $\beta \in \{0,1\}^{m_1}$ such that
$f(p_{\beta}) \in G \setminus Z(G)$ then there is an $a_i$ such that 
$[f(p_\beta),a_i] \neq 1$. Hence $g(p_\beta - i) \neq 1$ and 
$\val(\mathcal{J})\neq 1$ in $G \wr \Z$. 
This proves the theorem.
\end{proof}

\begin{thebibliography}{51}
\bibitem{ref1} Ian Agol. The virtual Haken conjecture. Documenta Mathematica, 18:1045–1087, 2013. With an appendix by Ian Agol, Daniel Groves, and Jason Manning.
\bibitem{ref2} Sanjeev Arora and Boaz Barak. Computational Complexity - A Modern Approach. Cambridge University Press, 2009. doi:10.1017/CBO9780511804090.
\bibitem{ref3} David A. Mix Barrington. Bounded-width polynomial-size branching programs recognize exactly those languages in $\mathrm{NC}^1$ . Journal of Computer and System Sciences, 38(1):150–164, 1989. doi:10.1016/0022-0000(89)90037-8.
\bibitem{ref4} David A. Mix Barrington and Denis Thérien. Finite monoids and the fine structure of $\mathrm{NC}^1$ . Journal of the ACM, 35:941–952, 1988. doi:10.1145/48014.63138.
\bibitem{ref5} Laurent Bartholdi, Michael Figelius, Markus Lohrey, and Armin Weiß. Groups with ALOGTIME-hard word problems and PSPACE-complete compressed word problems. Technical report, arXiv.org, 2020. arXiv:1909.13781.
\bibitem{ref6} Laurent Bartholdi, Rostislav I. Grigorchuk, and Zoran Šuniḱ. Branch groups. In Handbook of algebra, Volume 3, pages 989–1112. Elsevier/North-Holland, Amsterdam, 2003. doi: 10.1016/S1570-7954(03)80078-5.
\bibitem{ref7} Laurent Bartholdi and Volodymyr V. Nekrashevych. Iterated monodromy groups of quadratic polynomials. I. Groups, Geometry, and Dynamics, 2(3):309–336, 2008. doi:10.4171/GGD/42.
\bibitem{ref8} Martin Beaudry, Pierre McKenzie, Pierre Péladeau, and Denis Thérien. Finite monoids: From word to circuit evaluation. SIAM Journal on Computing, 26(1):138–152, 1997. doi: 10.1137/S0097539793249530.
\bibitem{ref9} Richard Beigel and John Gill. Counting classes: Thresholds, parity, mods, and fewness. Theoretical Computer Science, 103(1):3–23, 1992. doi:10.1016/0304-3975(92)90084-S.
\bibitem{ref10} William W. Boone. The Word Problem. Annals of Mathematics, 70(2):207–265, 1959. URL: www.jstor.org/stable/1970103.
\bibitem{ref11} Daniel P. Bovet, Pierluigi Crescenzi, and Riccardo Silvestri. A uniform approach to define complexity classes. Theoretical Computer Science, 104(2):263–283, 1992. doi:10.1016/ 0304-3975(92)90125-Y.
\bibitem{ref12} John W. Cannon, William J. Floyd, and Walter R. Parry. Introductory notes on Richard Thompson’s groups. L’Enseignement Mathématique, 42(3):215–256, 1996.
\bibitem{ref13} Moses Charikar, Eric Lehman, Ding Liu, Rina Panigrahy, Manoj Prabhakaran, Amit Sahai, and Abhi Shelat. The smallest grammar problem. IEEE Transactions on Information Theory, 51(7):2554–2576, 2005. doi:10.1109/TIT.2005.850116.
\bibitem{ref14} Max Dehn. Über unendliche diskontinuierliche Gruppen. Mathematische Annalen, 71(1):116– 144, 1911. doi:10.1007/BF01456932.
\bibitem{ref15} Michael R. Garey and David S. Johnson. Computers and Intractability: A Guide to the Theory of NP–completeness. Freeman, 1979.
\bibitem{ref16} Max Garzon and Yechezkel Zalcstein. The complexity of Grigorchuk groups with application to cryptography. Theoretical Computer Science, 88(1):83–98, 1991. doi:10.1016/0304-3975(91) 90074-C.
\bibitem{ref17} Rostislav I. Grigorchuk. Burnside’s problem on periodic groups. Functional Analysis and Its Applications, 14:41–43, 1980. doi:10.1007/BF01078416.
\bibitem{ref18} Rostislav I. Grigorchuk and Zoran Šuniḱ. Asymptotic aspects of Schreier graphs and Hanoi Towers groups. C. R. Math. Acad. Sci. Paris, 342(8):545–550, 2006. doi:10.1016/j.crma. 2006.02.001.
\bibitem{ref19} Victor S. Guba and Mark V. Sapir. On subgroups of the R. Thompson group F and other diagram groups. Matematicheskii Sbornik, 190(8):3–60, 1999. doi:10.1070/ SM1999v190n08ABEH000419.
\bibitem{ref20} Narain Gupta and Saïd Sidki. On the Burnside problem for periodic groups. Mathematische Zeitschrift, 182(3):385–388, 1983. doi:10.1007/BF01179757.
\bibitem{ref21} Frédéric Haglund and Daniel T. Wise. Coxeter groups are virtually special. Advances in Mathematics, 224(5):1890–1903, 2010. doi:10.1016/j.aim.2010.01.011.
\bibitem{ref22} Ulrich Hertrampf. Über Komplexitätsklassen, die mit Hilfe von k-wertigen Funktionen definiert werden. Habilitationsschrift, Universität Würzburg, 1994.
\bibitem{ref23} Ulrich Hertrampf. The shapes of trees. In Proceedings of COCOON 1997, volume 1276 of Lecture Notes in Computer Science, pages 412–421. Springer, 1997. doi:10.1007/BFb0045108.
\bibitem{ref24} Ulrich Hertrampf. Algebraic acceptance mechanisms for polynomial time machines. SIGACT News, 31(2):22–33, 2000. doi:10.1145/348210.348215.
\bibitem{ref25} Ulrich Hertrampf, Clemens Lautemann, Thomas Schwentick, Heribert Vollmer, and Klaus W. Wagner. On the power of polynomial time bit-reductions. In Proceedings of the Eighth Annual Structure in Complexity Theory Conference, pages 200–207. IEEE Computer Society Press, 1993. doi:10.1109/SCT.1993.336526.
\bibitem{ref26} Ulrich Hertrampf, Heribert Vollmer, and Klaus Wagner. On balanced versus unbalanced computation trees. Mathematical Systems Theory, 29(4):411–421, 1996. doi:10.1007/BF01192696.
\bibitem{ref27} Derek F. Holt, Markus Lohrey, and Saul Schleimer. Compressed decision problems in hyperbolic groups. In Proceedings of STACS 2019, volume 126 of LIPIcs, pages 37:1–37:16. Schloss Dagstuhl - Leibniz-Zentrum für Informatik, 2019. URL: http://www.dagstuhl.de/dagpub/ 978-3-95977-100-9.
\bibitem{ref28} Derek F. Holt, Sarah Rees, and Claas E. Röver. Groups, Languages and Automata, volume 88 of London Mathematical Society Student Texts. Cambridge University Press, 2017. doi: 10.1017/9781316588246.
\bibitem{ref29} Birgit Jenner, Pierre McKenzie, and Denis Thérien. Logspace and logtime leaf languages. Information and Computation, 129(1):21–33, 1996. doi:10.1006/inco.1996.0071.
\bibitem{ref30} Howard J. Karloff and Walter L. Ruzzo. The iterated mod problem. Information and Computation, 80(3):193–204, 1989. doi:10.1016/0890-5401(89)90008-4.
\bibitem{ref31} Daniel König and Markus Lohrey. Evaluation of circuits over nilpotent and polycyclic groups. Algorithmica, 80(5):1459–1492, 2018. doi:10.1007/s00453-017-0343-z.
\bibitem{ref32} Daniel König and Markus Lohrey. Parallel identity testing for skew circuits with big powers and applications. International Journal of Algebra and Computation, 28(6):979–1004, 2018. doi:10.1142/S0218196718500431.
\bibitem{ref33} Jörg Lehnert and Pascal Schweitzer. The co-word problem for the Higman-Thompson group is context-free. Bulletin of the London Mathematical Society, 39(2):235–241, February 2007. doi:10.1112/blms/bdl043.
\bibitem{ref34} Martin W. Liebeck, Eamonn A. O’Brien, Aner Shalev, and Pham Huu Tiep. The Ore conjecture. Journal of the European Mathematical Society, 12(4):939–1008, 2010. doi:10.4171/JEMS/220.
\bibitem{ref35} Yury Lifshits and Markus Lohrey. Querying and embedding compressed texts. In Proceedings of MFCS 2006, volume 4162 of Lecture Notes in Computer Science, pages 681–692. Springer, 2006. doi:10.1007/11821069\_59.
\bibitem{ref36} Richard J. Lipton and Yechezkel Zalcstein. Word problems solvable in logspace. Journal of the Association for Computing Machinery, 24(3):522–526, 1977. doi:10.1145/322017.322031.
\bibitem{ref37} Markus Lohrey. Leaf languages and string compression. Information and Computation, 209(6):951–965, 2011. doi:10.1016/j.ic.2011.01.009.
\bibitem{ref38} Markus Lohrey. The Compressed Word Problem for Groups. Springer Briefs in Mathematics. Springer, 2014. doi:10.1007/978-1-4939-0748-9.
\bibitem{ref39} Volodymyr Nekrashevych. Self-similar groups, volume 117 of Mathematical Surveys and Monographs. American Mathematical Society, Providence, RI, 2005. doi:10.1090/surv/117.
\bibitem{ref40} Piotr S. Novikov. On the algorithmic unsolvability of the word problem in group theory. Trudy Mat. Inst. Steklov, pages 1–143, 1955. In Russian. URL: http://mi.mathnet.ru/eng/tm1180.
\bibitem{ref41} David Robinson. Parallel Algorithms for Group Word Problems. PhD thesis, University of California, San Diego, 1993.
\bibitem{ref42} Joseph J. Rotman. An Introduction to the Theory of Groups (fourth edition). Springer, 1995.
\bibitem{ref43} Amir Shpilka and Amir Yehudayoff. Arithmetic circuits: A survey of recent results and open questions. Foundations and Trends in Theoretical Computer Science, 5(3-4):207–388, 2010. doi:10.1561/0400000039.
\bibitem{ref44} Hans-Ulrich Simon. Word problems for groups and contextfree recognition. In Proceedings of FCT 1979, pages 417–422. Akademie-Verlag, 1979.
\bibitem{ref45} Roman Smolensky. Algebraic methods in the theory of lower bounds for boolean circuit complexity. In Proceedings of STOC 1987, pages 77–82. ACM, 1987. doi:10.1145/28395. 28404.
\bibitem{ref46} Jacques Tits. Free subgroups in linear groups. Journal of Algebra, 20(2):250–270, 1972. doi:10.1016/0021-8693(72)90058-0.
\bibitem{ref47} Heribert Vollmer. Introduction to Circuit Complexity. Springer, Berlin, 1999. doi:10.1007/ 978-3-662-03927-4.
\bibitem{ref48} Stephan Waack. The parallel complexity of some constructions in combinatorial group theory. Journal of Information Processing and Cybernetics EIK, 26:265–281, 1990. doi: 10.1007/BFb0029647.
\bibitem{ref49} Jan Philipp Wächter and Armin Weiß. An automaton group with pspace-complete word problem. In Proceedings of STACS 2020, volume 154 of LIPIcs, pages 6:1–6:17. Schloss Dagstuhl - Leibniz-Zentrum für Informatik, 2020. URL: https://www.dagstuhl.de/dagpub/ 978-3-95977-140-5.
\bibitem{ref50} John S. Wilson. Embedding theorems for residually finite groups. Mathematische Zeitschrift, 174(2):149–157, 1980. doi:10.1007/BF01293535.
\bibitem{ref51} Daniel T. Wise. Research announcement: the structure of groups with a quasiconvex hierarchy. Electronic Research Announcements in Mathematical Sciences, 16:44–55, 2009. URL: http: //aimsciences.org//article/id/8d3fc128-8d02-4fee-af67-38001ca1d0ac.
\end{thebibliography}
\end{document}
